\section{Introdução}
\label{sec:introducao}
Filtros analógicos são circuitos eletrônicos essenciais no processamento de sinais, projetados para selecionar faixas de frequência específicas enquanto atenuam sinais interferentes indesejados \cite{razavi2006fund}. Em sistemas de comunicação sem fio e telefonia móvel, essa seletividade é crítica: receptores operam em ambientes congestionados e precisam isolar o canal de interesse de ruídos e interferências de bandas adjacentes antes que o sinal seja processado \cite{razavi2006fund}.

Diferentemente de aplicações em baixa frequência — como a filtragem de ruído da rede elétrica ($50\text{ Hz}$ ou $60\text{ Hz}$) em sinais de áudio —, o projeto de filtros em Radiofrequência (RF) exige que o fluxo do sinal seja tratado em termos de potência. Nessa faixa de operação, o comprimento de onda torna-se comparável às dimensões dos componentes, tornando indispensável o casamento de impedância (padronizado em $50~\Omega$) para garantir a máxima transferência de potência, minimizar perdas de inserção (\textit{insertion loss}) e evitar distorções de fase \cite{ludwig2000rf}.

Projetar filtros de alta frequência diretamente do zero torna-se matematicamente complexo à medida que a ordem aumenta, exigindo um método sistemático passo a passo \cite{ludwig2000rf}. Inicialmente, define-se a categoria de filtragem adequada à aplicação — passa-baixas (\textit{low-pass}), passa-altas (\textit{high-pass}), passa-faixa (\textit{bandpass}) ou rejeita-faixa (\textit{bandstop}) — partindo-se de um protótipo passa-baixas normalizado. Em seguida, seleciona-se a função de aproximação mais apropriada, como Butterworth ou Chebyshev, bem como a ordem do filtro ($N$), que determinam o compromisso entre a seletividade e o comportamento na banda passante, além da quantidade e distribuição dos elementos reativos. Por fim, os coeficientes normalizados são submetidos à desnormalização de impedância e à transformação em frequência para adequar os valores dos componentes à faixa de operação desejada \cite{ludwig2000rf}.

O relatório apresenta o dimensionamento e a validação computacional de um filtro passa-faixa de quinta ordem ($N = 5$) baseado na aproximação de Chebyshev, destinado a um receptor operando na frequência central de $f_0 = 6,5\text{ GHz}$ com largura de banda de $BW = 500\text{ MHz}$ (espectro de $6,25\text{ GHz}$ a $6,75\text{ GHz}$) e impedâncias terminais de $50~\Omega$, seguindo as especificações estabelecidas no roteiro da atividade \cite{Filter1}.

A escolha dessa faixa de operação fundamenta-se na sua relevância estratégica para os sistemas sem fio de alta capacidade e curto alcance. Por um lado, uma portadora em torno de $6,5\text{ GHz}$ com largura de banda de $500\text{ MHz}$ é característica de aplicações \textit{Ultra-Wideband} (UWB) — a exemplo do Canal 5 do padrão IEEE 802.15.4, centrado em aproximadamente $6,49\text{ GHz}$ —, nas quais o amplo espectro permite a transmissão de pulsos de curtíssima duração e elevada resolução temporal para localização precisa (ordem de centímetros), sensoriamento, radares de curto alcance e comunicação veicular (V2X) em veículos conectados \cite{ieee802154}. Nessa região de frequência, a maior atenuação por obstáculos favorece enlaces locais de alta taxa e precisão em detrimento de comunicações de longa distância.

Por outro lado, a banda de $6\text{ GHz}$ ($5,925\text{ GHz}$ a $7,125\text{ GHz}$) insere-se diretamente no atual cenário regulatório de compartilhamento espectral entre redes locais de alta velocidade e serviços móveis. No Brasil, a Agência Nacional de Telecomunicações (Anatel), por meio do Acórdão nº~203 e do Ato nº~10.400 de agosto de 2026, destinou oficialmente a subfaixa inferior de $5,925\text{ GHz}$ a $6,425\text{ GHz}$ ($500\text{ MHz}$ de largura) para equipamentos de radiação restrita (Wi-Fi 6E e Wi-Fi 7), mantendo os $700\text{ MHz}$ superiores ($6,425\text{ GHz}$ a $7,125\text{ GHz}$) alocados ao Serviço Móvel Pessoal (5G/6G), em consonância com a norma europeia ETSI EN 303 687 e as diretrizes da WRC-23 da UIT \cite{anatel2026ato10400, etsi303687}. Operar em $f_0 = 6,5\text{ GHz}$ posiciona o receptor exatamente nessa fronteira crítica de transição espectral, tornando rigoroso o requisito de atenuação mínima de $25\text{ dB}$ a partir de $1,5 \times BW$ para suprimir interferências de canais móveis adjacentes \cite{Filter1}.

